\documentclass[12pt,a4paper,oneside]{report}

\usepackage[main=slovak,english]{babel}
\usepackage{fontspec}
\usepackage{microtype}
\usepackage[a4paper,margin=25mm,headheight=15pt]{geometry}
\usepackage{setspace}
\onehalfspacing
\usepackage{amsmath}
\usepackage{graphicx}
\usepackage{float}
\usepackage{booktabs}
\usepackage{array}
\usepackage{tabularx}
\usepackage{enumitem}
\usepackage{xcolor}
\usepackage{tikz}
\usetikzlibrary{arrows.meta,babel,positioning}
\usepackage{csquotes}
\usepackage[hidelinks,unicode]{hyperref}

\definecolor{diagramblue}{HTML}{205781}
\definecolor{diagramlight}{HTML}{EAF3F8}
\definecolor{diagramgreen}{HTML}{E9F4ED}
\definecolor{diagramorange}{HTML}{FCF1DF}

\setlength{\parindent}{1.2cm}
\setlength{\parskip}{0.25em}
\setcounter{tocdepth}{2}
\setcounter{secnumdepth}{3}
\renewcommand{\arraystretch}{1.2}

\hypersetup{
  pdftitle={Semestrálna práca S1},
  pdfauthor={Bc. Pavol Pohánka},
  pdfsubject={Algoritmy a údajové štruktúry 2, akademický rok 2026/2027}
}

\begin{document}
\selectlanguage{slovak}

\begin{titlepage}
    \centering
    {\scshape\Large Žilinská univerzita v Žiline\par}
    \vspace{0.4cm}
    {\scshape\large Fakulta riadenia a informatiky\par}
    \vspace{1.2cm}
    {\large Algoritmy a údajové štruktúry 2\par}
    \vfill
    {\LARGE\bfseries Semestrálna práca S1\par}
    \vfill
    \noindent
    \begin{minipage}[t]{0.49\textwidth}
        \raggedright
        {\large Bc. Pavol Pohánka}
    \end{minipage}
    \begin{minipage}[t]{0.49\textwidth}
        \raggedleft
        {\large 2026/2027}
    \end{minipage}
\end{titlepage}

\pagenumbering{roman}
\tableofcontents
\clearpage
\pagenumbering{arabic}

\chapter{Architektúra aplikácie a návrh riešenia}

\section{Projekt a jeho časti}
Projekt vznikol v jazyku Java 21 a zostavuje sa nástrojom Maven.
Jeho základom sú tri vlastné priestorové stromy: KDTree, QuadTree
a QuadKDTree. Na porovnanie je pripojených šesť externých implementácií.
Súčasťou projektu je testovací nástroj a aplikácia POI Manager na
správu bodov záujmu. Aplikácia má grafické rozhranie JavaFX a obsahuje
aj terminálové rozhranie \texttt{ConsoleView}. Vstupným bodom je
trieda \texttt{MainApplication}.

\section{Architektúra a tok údajov}

POI Manager tvorí používateľské rozhranie, riadiaca vrstva a model
(obrázok \ref{fig:architecture}). \texttt{PoiManager\allowbreak Controller}
sprostredkúva požiadavky medzi rozhraním a \texttt{PoiManager\allowbreak Model}.
Model uchováva POI v zozname, na vyhľadávanie používa zvolený strom
a prácu so súbormi CSV zabezpečuje cez \texttt{Poi\allowbreak Csv\allowbreak Repository}.

\begin{figure}[H]
\centering
\begin{tikzpicture}[
    box/.style={
        draw=diagramblue,
        rounded corners=3pt,
        align=center,
        text width=4.6cm,
        minimum height=1cm,
        inner sep=4pt,
        fill=diagramlight,
        font=\footnotesize
    },
    small/.style={box, text width=3.6cm},
    edge/.style={-{Latex[length=2mm]}, draw=diagramblue, thick}
]

% Hlavný tok
\node[box] (view) at (0,0) {
    \textbf{\texttt{PoiManagerFxApplication}}\\
    {\scriptsize tabuľka, mapa, ovládacie prvky}
};
\node[box] (controller) at (0,-1.7) {
    \textbf{\texttt{PoiManagerController}}\\
    {\scriptsize sprostredkuje požiadavky}
};
\node[box] (model) at (0,-3.4) {
    \textbf{\texttt{PoiManagerModel}}\\
    {\scriptsize pois, tree, csvRepository}
};

% Pomocné triedy
\node[small, fill=diagramgreen] (poi) at (-4.6,-5.9) {
    \textbf{\texttt{Poi}}\\
    {\scriptsize údaje o bode záujmu}
};
\node[small, fill=diagramgreen] (tree) at (0,-5.9) {
    \textbf{\texttt{SpatialTree}}\\
    {\scriptsize priestorový index POI}
};
\node[small, fill=diagramorange] (csv) at (4.6,-5.9) {
    \textbf{\texttt{PoiCsvRepository}}\\
    {\scriptsize načítanie a zápis CSV}
};

% Hrany
\draw[edge] (view) -- (controller);
\draw[edge] (controller) -- (model);
\draw[edge] (model.south) -- (tree.north);
\draw[edge] (model.south) -- ++(0,-0.65) -| (poi.north);
\draw[edge] (model.south) -- ++(0,-0.65) -| (csv.north);

% Popisky hrán
\node[font=\scriptsize, above] at (-2.3,-4.55) {obsahuje 0..*};
\node[font=\scriptsize, above] at (2.3,-4.55) {používa};
\node[font=\scriptsize, right] at (0,-5.0) {používa};
\end{tikzpicture}
\caption{UML diagram hlavných tried POI Managera}
\label{fig:architecture}
\end{figure}

Grafické rozhranie vytvára \texttt{PoiManagerFxApplication}. Akcie
v \texttt{PoiControlsView} a udalosti v tabuľke či na mape spúšťajú
obslužné funkcie tejto triedy. Tie odovzdajú požiadavku triede
\texttt{PoiManager\allowbreak Controller}, ktorá zavolá príslušnú
operáciu modelu. Výsledky sa následne zobrazia
v \texttt{PoiTableView} a \texttt{PoiMapView}. Model nepracuje priamo
s JavaFX; používateľské rozhranie tak možno meniť bez zásahu do
priestorového vyhľadávania a práce s údajmi.

Balík \texttt{sk.fri.tester} obsahuje nástroj na overenie a porovnanie
stromov, \texttt{sk.fri.adapters} obsahuje adaptéry externých stromov.

\chapter{Priestorové dátové štruktúry}
\label{chap:structures}

\section{Spoločné rozhranie a body}
Priestorový strom organizuje objekty podľa polohy. Pri vyhľadávaní
v oblasti môže vynechať vetvy, ktoré neobsahujú hľadané body.

\texttt{SpatialTree<T,P>} je spoločná abstraktná trieda stromov,
kde \(T\) je typ objektu a \(P\) typ bodu. Poskytuje operácie:
\begin{itemize}[leftmargin=1.5em]
\item \texttt{insert(T)} vloží objekt a vráti informáciu o úspechu,
\item \texttt{find(P)} vyhľadá objekty na zadanej polohe,
\item \texttt{rangeSearch(P min, P max)} vyhľadá objekty v oblasti
  určenej dolnými a hornými súradnicami.
\end{itemize}
Pri vytvorení stromu sa odovzdáva funkcia na získanie súradníc objektu
a komparátor určujúci usporiadanie hodnôt súradníc.

Pre dvojrozmerné údaje sa používa \texttt{Point<K>} s dvojicou
súradníc a pre \(k\)-rozmerný KDTree slúži \texttt{KDPoint<K>} so
zoznamom súradníc. Pri hľadaní bodu sa musia zhodovať všetky jeho
súradnice podľa komparátora.

Objekty s rovnakými súradnicami sa ukladajú do zoznamu hodnôt jedného
uzla. Pri vyhľadávaní v oblasti sa započítajú aj body na jej okraji.

Vkladanie a hľadanie používajú cykly, pri prehľadávaní viacerých vetiev
sa používa zásobník.

\section{KDTree}
\texttt{KDTree<T,K>} rozdeľuje priestor postupne podľa jednej súradnice.
V dvojrozmernom prípade sa striedajú súradnice \(x\) a \(y\), všeobecne
sa používa dimenzia \(d \bmod k\), kde \(d\) je hĺbka uzla.
\texttt{KDNode} obsahuje bod, jeho hodnoty a odkazy na dvoch potomkov.
Pri vkladaní a hľadaní bodu sa pokračuje vetvou podľa deliacej súradnice,
pri jej zhode a rozdielnom bode sa pokračuje vpravo.

Príklad delenia je na obrázku \ref{fig:kd-partition}.

\begin{figure}[H]
\centering
\begin{tikzpicture}[font=\small]
\fill[diagramlight] (-4,-2.3) rectangle (4,2.3);
\draw[thick, diagramblue] (-4,-2.3) rectangle (4,2.3);
\draw[very thick, diagramblue] (0,-2.3) -- (0,2.3);
\draw[very thick, diagramblue, dashed] (-4,1) -- (0,1);
\draw[very thick, diagramblue, dashed] (0,-0.8) -- (4,-0.8);
\fill[red!75!black] (0,0) circle (2.5pt);
\fill[red!75!black] (-2.2,1) circle (2.5pt);
\fill[red!75!black] (2,-0.8) circle (2.5pt);
\node[below right] at (0,0) {1. bod};
\node[above left] at (-2.2,1) {2. bod};
\node[above right] at (2,-0.8) {3. bod};
\node[above] at (0,2.3) {delenie podľa \(x\)};
\node[below] at (-2,-2.3) {delenie podľa \(y\)};
\node[below] at (2,-2.3) {delenie podľa \(y\)};
\end{tikzpicture}
\caption{Príklad striedania súradníc pri delení v KDTree}
\label{fig:kd-partition}
\end{figure}

Pri vyhľadávaní v oblasti strom kontroluje navštívené body a podľa
hraníc oblasti vyberá vetvy, v ktorých sa môže nachádzať výsledok.
Príklad je na obrázku \ref{fig:range-search}.

\begin{figure}[H]
\centering
\begin{tikzpicture}[font=\small]
\fill[gray!12] (-4,-2) rectangle (0,2);
\fill[diagramlight] (0,-2) rectangle (4,2);
\draw[thick, diagramblue] (-4,-2) rectangle (4,2);
\draw[very thick, diagramblue] (0,-2) -- (0,2);
\fill[diagramgreen] (1,-0.6) rectangle (3.4,1.3);
\draw[thick, green!45!black] (1,-0.6) rectangle (3.4,1.3);
\fill[red!75!black] (0,0) circle (2.5pt);
\fill[red!75!black] (1.7,0.5) circle (2.5pt);
\fill[red!75!black] (2.8,-1.2) circle (2.5pt);
\node at (-2,0) {vynechaná vetva};
\node[above] at (0,2) {deliaca súradnica};
\node[above] at (2.2,1.3) {hľadaná oblasť};
\end{tikzpicture}
\caption{Príklad vyhľadávania v oblasti ležiacej napravo od delenia}
\label{fig:range-search}
\end{figure}

Pri výške stromu \(h\) trvá vkladanie a hľadanie bodu \(O(hk)\).
Bez vyvažovania môže byť \(h\) až \(n\), počet rôznych bodov.
Vyhľadávanie v oblasti trvá \(O(vk+r)\), kde \(v\) je počet navštívených
uzlov a \(r\) počet výsledkov, v najhoršom prípade \(O(nk+r)\).
Pre pevné \(k\) je pamäťová náročnosť lineárna vzhľadom na počet uzlov
a uložených hodnôt.

\section{QuadTree}
\texttt{QuadTree<T,K>} rozdeľuje rovinu na štyri kvadranty podľa bodu
uzla (obrázok \ref{fig:quadrants}). Uzol obsahuje bod, jeho hodnoty
a až štyroch potomkov. Pri vkladaní aj hľadaní bodu sa podľa súradníc
vyberie príslušný kvadrant. Body na deliacej priamke patria napravo
alebo nahor.

\begin{figure}[H]
\centering
\begin{tikzpicture}[scale=0.82]
\fill[diagramlight] (-4,-3) rectangle (4,3);
\draw[thick, diagramblue] (-4,-3) rectangle (4,3);
\draw[dashed,diagramblue] (0,-3) -- (0,3);
\draw[dashed,diagramblue] (-4,0) -- (4,0);
\fill[red!80!black] (0,0) circle (2.5pt);
\node[above left] at (-0.1,0.1) {\small uzol \((x_0,y_0)\)};
\node at (-2,1.5) {\small \(x<x_0,\ y\geq y_0\)};
\node at (2,1.5) {\small \(x\geq x_0,\ y\geq y_0\)};
\node at (-2,-1.5) {\small \(x<x_0,\ y<y_0\)};
\node at (2,-1.5) {\small \(x\geq x_0,\ y<y_0\)};
\draw[-{Latex[length=2mm]}, thick] (0,-3.35) -- (3.7,-3.35) node[right] {\small \(x\)};
\draw[-{Latex[length=2mm]}, thick] (-4.35,0) -- (-4.35,2.7) node[above] {\small \(y\)};
\end{tikzpicture}
\caption{Rozdelenie do kvadrantov podľa súradníc uzla QuadTree}
\label{fig:quadrants}
\end{figure}

Pri vyhľadávaní v oblasti strom prechádza kvadranty, ktoré sa môžu
prekrývať so zadanou oblasťou. Vkladanie a hľadanie bodu trvá \(O(h)\),
vyhľadávanie v oblasti \(O(v+r)\), najviac \(O(n+r)\). Pamäťová
náročnosť je lineárna vzhľadom na počet uzlov a hodnôt.

Maximálna hĺbka je 32. Nový bod vyžadujúci hlbší uzol sa odmietne,
k existujúcemu bodu možno pridať ďalšiu hodnotu.

\section{QuadKDTree}
\texttt{QuadKDTree<T,K>} kombinuje kvadrantové delenie s KDTree.
Vonkajšia časť rozdeľuje priestor na kvadranty, zatiaľ čo body v každom
liste sú uložené v samostatnom KDTree. Hodnoty s rovnakými súradnicami
sa ukladajú spolu v uzle tohto KDTree.

List sa začne deliť pri vkladaní nového bodu, keď už obsahuje 32 bodov
s rôznymi súradnicami a ešte nedosiahol limit hĺbky. Jeho hodnoty sa
podľa polohy rozdelia do nových listov, ako ukazuje obrázok
\ref{fig:quadkd-split}. Kvadrantové delenie má maximálnu hĺbku 32.
Po dosiahnutí tohto limitu sa ďalšie body pridávajú do KDTree v liste,
ktorý sa už ďalej nedelí.

\begin{figure}[H]
\centering
\begin{tikzpicture}[
  box/.style={draw=diagramblue, rounded corners=3pt, align=center,
    minimum height=0.9cm, inner sep=4pt, font=\small},
  edge/.style={-{Latex[length=2mm]}, draw=diagramblue, thick}
]
\node[box, fill=diagramgreen, text width=2.8cm] (before) at (-4.2,0)
  {List s KDTree\\32 rôznych bodov};
\node[box, fill=diagramlight, text width=3cm] (root) at (2.2,0.6)
  {Kvadrantové\\rozdelenie};
\node[box, fill=diagramgreen, text width=1.5cm] (one) at (-1.1,-1.7) {KDTree\\v liste};
\node[box, fill=diagramgreen, text width=1.5cm] (two) at (1.1,-1.7) {KDTree\\v liste};
\node[box, fill=diagramgreen, text width=1.5cm] (three) at (3.3,-1.7) {KDTree\\v liste};
\node[box, fill=diagramgreen, text width=1.5cm] (four) at (5.5,-1.7) {KDTree\\v liste};
\draw[edge] (before) -- node[above, align=center, font=\scriptsize]
  {ďalší\\nový bod} (root);
\draw[edge] (root) -- (one);
\draw[edge] (root) -- (two);
\draw[edge] (root) -- (three);
\draw[edge] (root) -- (four);
\end{tikzpicture}
\caption{Príklad rozdelenia listu QuadKDTree do štyroch obsadených kvadrantov}
\label{fig:quadkd-split}
\end{figure}

Pri hľadaní bodu strom nájde príslušný list a prehľadá jeho KDTree.
Pri vyhľadávaní v oblasti prehľadá KDTree v príslušných kvadrantoch.
Pri výške vonkajšej cesty \(h\) a výške KDTree v liste \(h_k\) trvá
hľadanie bodu aj vkladanie bez delenia \(O(h+h_k k)\). Pri delení treba
navyše rozmiestniť všetky hodnoty listu, vrátane hodnôt s rovnakými
súradnicami. Vyhľadávanie v oblasti trvá v najhoršom prípade \(O(nk+r)\).
Pre pevné \(k\) je pamäťová náročnosť lineárna vzhľadom na počet uzlov
a uložených hodnôt.

\section{Porovnanie vlastných implementácií}
\begin{table}[H]
\centering
\small
\begin{tabularx}{\textwidth}{@{}p{2.6cm}XXX@{}}
\toprule
\textbf{Vlastnosť} & \textbf{KDTree} & \textbf{QuadTree} & \textbf{QuadKDTree}\\
\midrule
Rozdelenie priestoru & Striedanie súradnicových dimenzií & Kvadranty podľa bodu uzla & Kvadranty, KDTree v listoch\\
Dimenzie & \(k\) & 2 & 2\\
Duplicitné body & Zoznam hodnôt v uzle & Zoznam hodnôt v uzle & Zoznam hodnôt v uzle KDTree\\
Vyhľadávanie v oblasti & Výber vetiev podľa deliacej súradnice & Výber kvadrantov podľa oblasti & Výber kvadrantov a vyhľadávanie v KDTree listov\\
Limity & -- & Hĺbka 32, nové body nad limit sa odmietajú & Hĺbka 32, delenie listu pri 33. rôznom bode\\
\bottomrule
\end{tabularx}
\caption{Vlastnosti vlastných priestorových štruktúr}
\label{tab:own-trees}
\end{table}

\chapter{Externé implementácie}
\label{chap:external}

Externé stromy umožňujú porovnať vlastné riešenie s ďalšími
implementáciami. Každý z nich má adaptér odvodený od \texttt{SpatialTree}.
Adaptér prevádza body a výsledky do formátu projektu, aby testovací
nástroj mohol stromy používať rovnakým spôsobom.

\section{Externé KDTree}
\begin{itemize}[leftmargin=1.5em]
\item \textbf{KDtree-AlgorithmsAndStructures} \cite{externalKd} používa
  adaptér \texttt{KDAdapter}. Vkladanie a hľadanie bodu vykonáva strom,
  vyhľadávanie v oblasti adaptér rieši lineárnym prechodom cez pomocný
  zoznam bodov.
\item \textbf{TinSpin Indexes} \cite{tinspin} sprístupňuje
  \texttt{TinSpinKDAdapter}. Súradnice prevádza na pole hodnôt
  \texttt{double} a vyhľadávanie vykonáva externý strom.
\item \textbf{JTS} \cite{jts} používa \texttt{JtsKDAdapter}.
  Pre objekty s rovnakou polohou vedie pomocnú mapu hodnôt, oblasti
  zadáva pomocou triedy \texttt{Envelope}.
\end{itemize}

\section{Externé QuadTree}
\begin{itemize}[leftmargin=1.5em]
\item \textbf{Varun Pant Quadtree} \cite{varunQuad} používa
  \texttt{VarunPantQuadTreeAdapter}, ktorý zoskupuje hodnoty na rovnakej
  polohe. Rozsah stromu je nastavený na zemepisné súradnice.
\item \textbf{Alwex QuadTree} \cite{alwexQuad} sprístupňuje
  \texttt{AlwexQuadTreeAdapter}. Bod vkladá ako malý obdĺžnik a výsledné
  objekty filtruje podľa pôvodných súradníc. Implementácia používa
  súradnice typu \texttt{float}.
\item \textbf{Java Simple Quadtree} \cite{ryanQuad} používa
  \texttt{RyanQuadTreeAdapter}. Body vkladá ako obdĺžniky nulovej veľkosti
  a výsledky filtruje podľa zadanej oblasti vrátane jej hraníc.
\end{itemize}

\chapter{Testovanie a porovnanie výkonnosti}
\label{chap:testing}

\section{Testovací nástroj a dáta}
\texttt{TreeTester} overuje tri vlastné stromy a šesť externých adaptérov.
Vložené body vyhľadáva podľa polohy a výsledky hľadania v oblasti
porovnáva s lineárnym prechodom prijatých bodov. Zo súborov
\texttt{slovakia\_nodes.csv} a \texttt{testovacie\_oblasti.csv}
načíta najviac 10\,000 bodov a 1\,000 oblastí.

\section{Kontrola správnosti}
Všetkých deväť konfigurácií prešlo kontrolami vyhľadávania.
Pri kontrole vkladania a hľadania polohy sa vyberalo
1\,000 bodov s opakovaním, pri hľadaní v oblasti sa použilo 10\,000
bodov a 1\,000 oblastí. Vlastný QuadTree odmietol 4\,900 bodov pre limit
hĺbky 32, preto sa výsledky porovnávali iba s prijatými bodmi.

\section{Metodika benchmarku a obmedzenia}
\texttt{TreeBenchmark} meria vkladanie, hľadanie vložených bodov
a vyhľadávanie v oblasti. Výsledkom je medián časov v milisekundách.
Každý beh vytvára nový strom, počet hľadaní sa môže líšiť podľa počtu
úspešne vložených bodov.

Benchmark vykoná dva zahrievacie a päť meraných behov. Nástroj umožňuje
použiť až 1\,000 oblastí na beh, v uvedenom meraní bolo použitých 100.

\section{Výsledky testov}
Tabuľka uvádza mediány piatich behov s 10\,000 vloženiami,
10\,000 hľadaniami bodov a 100 dopytmi na oblasť na beh. Meranie sa
uskutočnilo pred opravou počtu zahrievacích behov, po oprave sa benchmark
znovu nespúšťal.

\begin{table}[H]
\centering
\small
\begin{tabularx}{\textwidth}{@{}Xrrrr@{}}
\toprule
\textbf{Implementácia} & \textbf{Vkladanie} & \textbf{Hľadanie} &
\textbf{Oblasť} & \textbf{Spolu}\\
\midrule
My KDTree & 5,266 & 5,056 & 0,512 & 10,833\\
My QuadTree & 5,845 & 6,182 & 0,806 & 12,833\\
My Quad-KDTree & 5,069 & 4,694 & 2,913 & 12,676\\
Existing external KDTree & 6,242 & 6,264 & 21,801 & 34,307\\
TinSpin KDTree & 4,254 & 25,984 & 0,446 & 30,683\\
JTS KDTree & 11,958 & 13,796 & 0,325 & 26,078\\
Varun Pant QuadTree & 4,171 & 4,316 & 0,069 & 8,557\\
Alwex QuadTree & 18,314 & 128,441 & 5,067 & 151,822\\
Ryan Pelletier QuadTree & 3,069 & 1\,139,936 & 10,828 & 1\,153,832\\
\bottomrule
\end{tabularx}
\caption{Výsledky benchmarku: mediány v milisekundách}
\label{tab:historical-results}
\end{table}

Najnižší celkový čas v uvedenom meraní má Varun Pant Quadtree
(8,557 ms). Spomedzi vlastných štruktúr dosiahol najnižší celkový čas
KDTree (10,833 ms).

\chapter{Aplikácia POI Manager}
\label{chap:poi}

\section{Záznamy a priestorový index}
\texttt{Poi} obsahuje identifikátor, názov, opis, súradnice, adresu
a čas vytvorenia a poslednej úpravy. Model ukladá záznamy do zoznamu
a ich polohu indexuje zvoleným stromom. Súradnice sú v poradí
\texttt{(latitude, longitude)}.

\section{Vyhľadávanie a správa POI}
Vyhľadávanie podľa polohy vráti aktívne POI na zadaných súradniciach.
Pri hľadaní v okolí sa najprv vyberú body v obdĺžniku a potom sa overí
ich vzdialenosť. Pri hľadaní názvu v oblasti sa výsledky filtrujú podľa
názvu bez ohľadu na veľkosť písmen a diakritiku.

POI možno vytvárať, upravovať, zobrazovať a mazať. Mazanie je logické:
záznam zostáva v pamäti, ale bežné vyhľadávanie ho vynechá. Pri úprave
sa index zostaví znova. Skúšobné POI sa generujú v obdĺžnikovej oblasti,
preto môžu ležať aj mimo hranice Slovenska.

\section{CSV a používateľské rozhranie}
CSV súbor obsahuje iba aktívne záznamy. Polia sú oddelené bodkočiarkou,
preto názov, opis ani adresa nesmú obsahovať bodkočiarku alebo nový riadok.
Zmeny sa ukladajú akciou používateľa.

Rozhranie JavaFX zobrazuje POI v tabuľke a na offline mape Slovenska
\cite{naturalEarth}. Mapa umožňuje posun, zmenu mierky, výber bodu
a vytvorenie POI kliknutím.

\chapter{Používateľská príručka}

\section{Spustenie}
Na spustenie je potrebný JDK 21 a Maven. V koreňovom adresári projektu
spustite:
\begin{verbatim}
mvn javafx:run
\end{verbatim}
K dispozícii je aj spustiteľná verzia aplikácie pre Windows vo formáte
\texttt{.exe}.
Po spustení vyberte v úvodnom dialógu KD Tree, Quad Tree alebo Quad KD
Tree. V nasledujúcom kroku zvoľte \enquote{New file} a určte nový CSV
súbor, alebo vyberte \enquote{Open file} a načítajte existujúci CSV
súbor s bodmi záujmu. Zvolený strom zostane aktívny počas celého behu aplikácie,
zmeniť ho možno až po jej opätovnom spustení.

Projekt obsahuje aj terminálové rozhranie \texttt{Console\allowbreak View},
jeho spustenie je v \texttt{Main\allowbreak Application} zakomentované.

\section{Každodenná práca}
\begin{enumerate}[leftmargin=1.5em]
\item Na hľadanie podľa polohy zvoľte \enquote{Find by GPS} a zadajte súradnice v desatinných stupňoch.
\item Na vyhľadanie okolia kliknite na riadok POI v tabuľke, použite \enquote{Find nearby} a zadajte polomer v metroch. Stredom je poloha vybraného záznamu.
\item Na hľadanie podľa názvu použite \enquote{Search by name} a zadajte text aj hranice oblasti.
\item Na vytvorenie POI použite \enquote{Add POI}, prípadne kliknite do mapy vnútri hranice Slovenska. Na úpravu zaškrtnite práve jeden aktívny záznam a zvoľte \enquote{Edit selected}.
\item Na vymazanie zaškrtnite záznamy a použite \enquote{Delete selected}. Cez \enquote{Show deleted} možno zobraziť aj vymazané záznamy.
\item Rozsah zobrazených záznamov sa mení v ponuke \enquote{Show}, \enquote{All POI} zobrazí všetky. \enquote{Generate test POIs} vytvorí skúšobné záznamy.
\item Pred zatvorením alebo prepnutím súboru uložte zmeny cez \enquote{Save CSV}, automatické ukladanie nie je zapnuté.
\end{enumerate}

\chapter{Záver}

Výsledkom práce sú tri vlastné priestorové štruktúry, šesť externých
adaptérov a aplikácia POI Manager, ktorá index využíva na hľadanie POI
podľa polohy, oblasti a vzdialenosti. V zaznamenanom teste správnosti
deväť konfigurácií prešlo kontrolami nad skúšobnou vzorkou, vlastný
QuadTree však pre limit hĺbky odmietol časť bodov. Benchmarkové výsledky
zaznamenávajú časy namerané pri použitom nastavení testu.

\renewcommand{\bibname}{Zoznam použitej literatúry}
\begin{thebibliography}{10}
\addcontentsline{toc}{chapter}{Zoznam použitej literatúry}

\bibitem{externalKd} Luka3691.
\textit{KDtree-AlgorithmsAndStructures}.
\url{https://github.com/luka3691/KDtree-AlgorithmsAndStructures}.
[cit. 2026-11-09].

\bibitem{tinspin} Tzaeschke.
\textit{TinSpin Indexes}.
\url{https://github.com/tzaeschke/tinspin-indexes}.
[cit. 2026-11-09].

\bibitem{jts} LocationTech.
\textit{JTS Topology Suite}.
\url{https://github.com/locationtech/jts}.
[cit. 2026-11-09].

\bibitem{varunQuad} Varun Pant.
\textit{Quadtree}.
\url{https://github.com/varunpant/Quadtree}.
[cit. 2026-11-09].

\bibitem{alwexQuad} Alwex.
\textit{QuadTree}.
\url{https://github.com/alwex/QuadTree}.
[cit. 2026-11-09].

\bibitem{ryanQuad} Ryan Pelletier.
\textit{Java Simple Quadtree}.
\url{https://github.com/ryanmpelletier/java-simple-quadtree}.
[cit. 2026-11-09].

\bibitem{naturalEarth} Natural Earth.
\textit{Natural Earth -- public domain map data}.
\url{https://www.naturalearthdata.com/}.
[cit. 2026-11-09]. Zdrojová poznámka je uvedená v
\texttt{src/main/resources/map/slovakia-border.txt}.
\end{thebibliography}

\end{document}